\section{Results}
\label{sec:results}

Results are reported in the same H1--H4 order as the experimental design. All
200 planned prediction-data rollouts completed. Raw and merged manifests agree
on 11,230 samples; the 160/20/20 rollout split has disjoint groups, no duplicate
keys and no normalisation leakage. The largest descriptive split-covariate
standardised mean difference is 0.208. The formal closed-loop experiment also
completed all 80/80 prespecified rollouts and passed its treatment, integrity,
stopping and archive checks.

\subsection{H1: task adaptation improves in-distribution prediction}

% EVIDENCE: H1_B1_MINUS_B0_TEST_NLL, H1_B1_MINUS_B0_TEST_ADE,
% H1_B1_MINUS_B0_TEST_FDE
% SOURCE: supervisor_feedback_v1/03_finetune_audit/FINETUNE_AUDIT_MANIFEST.json
Across the three validation seeds, median rollout-macro NLL ranks \Bone{}
first, followed by \BtwoD{}, \Ttwo{}, \Tone{} and \BtwoM{}. On the one-shot
frozen test, \Bone{} reduces NLL from 2.171 to 1.857~nats/step, top-1 ADE from
1.283 to 0.100~m and FDE from 2.644 to 0.121~m relative to pretrained \Bzero{}.
The corresponding rollout-macro differences are $-0.314$~nats/step,
$-1.183$~m and $-2.523$~m. All three metrics favour \Bone{} in every one of
the five held-out initialisation groups; at that independent-unit level, the
mean paired effects are $-0.316$~nats/step NLL, $-1.189$~m ADE and
$-2.544$~m FDE. The two-sided exact sign-flip value reaches its attainable
minimum of $p=0.0625$, so the evidence is directionally consistent but is not
reported as conventional $p<0.05$ significance. \Bone{} also improves ADE and
FDE relative to constant velocity, clipped constant acceleration and the
train-mean prior in all five groups. H1 is therefore supported for the frozen
Town05 give-way distribution.

\begin{table}[t]
\centering\small
\caption{Frozen B0/B1 test results at one common rollout-macro aggregation. Every metric is first averaged within rollout and then over 20 rollouts. The 315 overlapping windows are descriptive observations, not independent replications; lower is better.}
\label{tab:finetune-b0-b1-rollout-macro}
\begin{tabular}{@{}lrrrrrr@{}}
\toprule
Stack/contrast & Windows & Rollouts & Init groups & NLL & ADE (m) & FDE (m) \\
\midrule
B0 & 315 & 20 & 5 & 2.171 & 1.283 & 2.644 \\
B1 & 315 & 20 & 5 & 1.857 & 0.100 & 0.121 \\
\midrule
B1$-$B0 & 315 & 20 & 5 & -0.314 & -1.183 & -2.523 \\
\bottomrule
\end{tabular}
\end{table}


The mechanism checks narrow this conclusion. Raster shuffling increases
\Bone{} ADE by approximately 0.270--0.293~m across seeds, whereas target-history
shuffling changes aggregate ADE by about 0.00008~m on average; the aggregate
gain is therefore raster-dominant. In a small response-active diagnostic,
uncalibrated \Bone{} NLL is 2.076 versus 2.472 for \Bzero{}, but the
validation-frozen calibration changes these values to 8.573 and 2.959. Because
this subset contains only 15 overlapping windows from six rollouts and three
initialisation groups, it identifies a possible calibration boundary rather
than a confirmatory failure rate.

An earlier progress report described an improvement from $0.98\%$ to $100\%$.
That percentage measured how often the top-probability mode matched the
oracle-best mode; it was neither a trajectory-accuracy threshold nor a valid
independent-sample estimate. It is therefore withdrawn. The frozen-test
NLL/ADE/FDE comparison above, using common rollout-macro aggregation and five
paired held-out initialisations, is the evidence used for H1.

\begin{center}
  \refstepcounter{figure}
  \label{fig:offline}
  \centering
  \includegraphics[width=0.98\linewidth]{figure03_offline_model_comparison.png}
  \smallskip
  {\small Figure~\thefigure: Frozen offline model comparison. Points show validation seeds and
  horizontal segments their medians; test ADE/FDE include the pretrained
  \Bzero{} control. NLL, ADE and FDE use rollout-macro aggregation, so
  overlapping windows are not counted as independent replications.}
\end{center}

\subsection{H2: the tested Transformer adapters add no consistent value}

% EVIDENCE: H2_T1_TEST_NLL, H2_B2_M_TEST_NLL,
% H2_T2_TEST_NLL, H2_B2_D_TEST_NLL
% SOURCE: distinction_v1/10_four_hypothesis_evidence/M1_EVIDENCE_MANIFEST.json
The two matched-scope comparisons point in opposite directions. For the
mean-only scope, \Tone{} reaches a frozen-test NLL of 2.003746 compared with
2.024409 for the MLP control \BtwoM{}, an improvement of 0.020663. For the
distributional scope, \Ttwo{} reaches 1.878148 compared with 1.872789 for
\BtwoD{}, a degradation of 0.005359. Neither Transformer outperforms the
simple \Bone{} result of 1.857094. H2 is therefore not supported under the
tested architecture and optimisation budget.

Zeroing or shuffling the valid interaction sequence worsens both Transformer
variants, showing that they use the temporal input; sequence neglect does not
explain the negative result. The correct interpretation is nevertheless
bounded. \Bone{} has approximately 1.034 million trainable parameters, while
the adapters contain approximately 77--176 thousand, and 10/15 training runs
select a checkpoint at the 20-epoch ceiling. The experiment therefore supports
simple adaptation as the strongest tested configuration, not a general claim
that Transformers cannot improve interaction prediction.

\subsection{H3: offline gain does not transfer consistently to closed loop}

The manipulation check succeeds before closed-loop outcomes are interpreted.
Across 40/40 prespecified metric--policy--style checks, in-loop predictions
favour \Bone{} over \Bzero{} in all five paired initialisation groups. The
\Bone{}--\Bzero{} top-1 ADE effects range from $-1.352$ to $-1.125$~m and FDE
effects from $-3.152$ to $-2.315$~m. Thus the weak translation cannot be
explained by the two deployed predictor stacks becoming indistinguishable.

Only two of the eight risk-policy--target-style cells jointly show faster mean
completion and no-worse mean physical separation for \Bone{} relative to
\Bzero{}: adaptive/reactive and fixed-aggressive/assertive. The remaining six
fail at least one registered direction. For example, adaptive/assertive is
0.240~s slower with 0.0067~m lower separation, whereas fixed-aggressive/
assertive is 0.790~s faster with 0.0053~m higher separation. All Holm-adjusted
exact tests are non-confirmatory. H3 is therefore not supported as a consistent
offline-to-closed-loop translation claim, although the two favourable cells
show that the adapted stack can have context-specific closed-loop value.

\begin{center}
  \refstepcounter{table}
  \label{tab:h3-translation}
  \small
  Table~\thetable: Predictor-stack translation in the formal closed-loop experiment.
  A supported direction requires lower mean completion time and no-worse mean
  separation for \Bone{}--\Bzero{}. Each cell contains five paired
  initialisation groups.\par\smallskip
  \begin{tabular}{@{}llrrl@{}}
    \toprule
    Policy & Style & Completion (s) & Separation (m) & Direction \\
    \midrule
    Adaptive & Assertive &  0.240 & $-0.0067$ & Not supported \\
    Adaptive & Reactive & $-0.180$ &  0.0006 & Supported \\
    Fixed aggressive & Assertive & $-0.790$ &  0.0053 & Supported \\
    Fixed aggressive & Reactive & $-0.010$ & $-0.0106$ & Not supported \\
    Fixed conservative & Assertive & $-0.520$ & $-0.0262$ & Not supported \\
    Fixed conservative & Reactive &  0.490 & $-0.0056$ & Not supported \\
    Fixed medium & Assertive &  0.230 &  0.0072 & Not supported \\
    Fixed medium & Reactive & $-0.220$ & $-0.0084$ & Not supported \\
    \bottomrule
\end{tabular}
\end{center}

Figure~\ref{fig:h3-translation} visualises the registered joint-direction
criterion: a favourable cell must improve completion without sacrificing the
measured physical margin, rather than improving either outcome alone.

\begin{center}
  \refstepcounter{figure}
  \label{fig:h3-translation}
  \centering
  \includegraphics[width=0.99\linewidth]{figure_r3_h3_translation.png}
  \smallskip
  {\small Figure~\thefigure: Predictor-stack translation across risk policies
  and target styles. All values are derived from the formal 80-rollout
  closed-loop analysis.}
\end{center}

All adaptive cells use the adaptive solver on every audited step, every one of
the 20 adaptive rollouts varies its risk allocation within rollout, and the
fixed settings remain stable across initialisations. All 80 rollouts complete,
with no observed native collision, zero-margin footprint collision,
fixed-geometry yield failure or completion failure. These zero counts establish
nominal reliability in the tested sample; they do not prove zero population
risk, mathematical MPC feasibility or equivalence between treatments.

\subsection{H4: adaptive risk is context dependent rather than dominant}

% EVIDENCE: H4_DOMINANCE_CELLS, H4_PRESPECIFIED_CELLS
% SOURCE: distinction_v1/10_four_hypothesis_evidence/M1_EVIDENCE_MANIFEST.json
Adaptive risk satisfies the prespecified empirical dominance rule in only
3/12 predictor--style--fixed-comparator contrasts: \Bzero{}/assertive versus
fixed medium, and \Bone{}/reactive versus fixed aggressive and fixed medium.
It does not dominate in the remaining 9/12 contrasts. H4 is therefore not
supported as a universal frontier-dominance claim. The positive comparisons
retain a narrower conclusion: adaptive risk can be a useful context-dependent
operating point, but it is not a general replacement for fixed risk.

\begin{center}
  \refstepcounter{figure}
  \label{fig:h4-dominance}
  \centering
  \includegraphics[width=0.99\linewidth]{figure_r3_h4_dominance.png}
  \smallskip
  {\small Figure~\thefigure: Adaptive-risk dominance across predictor stack,
  target style and fixed comparator. Dominance is the prespecified empirical
  Pareto rule with binary guards, not an equivalence test.}
\end{center}

Solver and supervisor telemetry limits the mechanism claim. Historical
logger-unaccepted rows comprise bypass/no-solve and attempted-but-nonaccepted
contexts and cannot be relabelled wholesale as infeasible MPC optimisations.
Finite CasADi solve-time and controller-acceptance rates therefore use factual
solve attempts as their denominator, while bypass requests, raw return status
and rollout outcome remain separate. Earlier timing diagnostics show that
physical margin, controller acceptance/fallback and supervisor activity can
change together; they do not identify the supervisor as the sole cause of the
H4 pattern.

% Provisional tracked insertion. The final-only evidence builder atomically
% replaces this file after SF1--SF4 closure receipts and their source hashes pass.
\subsection{Supervisor-feedback closure analyses}
\label{sec:supervisor-feedback-closure-results}

The final behaviour, finite solver-time, controller-acceptance, downstream-
outcome and corrected-supervisor application-authority results are intentionally unavailable
in this non-submission draft.  The release audit fails closed until the immutable
R3 raw-log analyses and the preregistered SF4 matrix have passed their scientific
and provenance gates.


\begin{samepage}
\begin{center}
  \refstepcounter{table}
  \label{tab:hypotheses}
  \small
  Table~\thetable: Evidence-bounded verdicts for the four headline
  hypotheses. Counts are prespecified cells or comparisons, not simulation
  steps.\par\smallskip
  \begin{tabular}{@{}L{0.06\linewidth}L{0.46\linewidth}L{0.20\linewidth}L{0.18\linewidth}@{}}
    \toprule
    ID & Hypothesis & Verdict & Boundary \\
    \midrule
    H1 & Task adaptation improves in-distribution prediction & Supported with boundary & Frozen Town05 distribution \\
    H2 & Transformer adapters add consistent value & Not supported & Tested configurations/budget \\
    H3 & Offline gain transfers consistently to closed loop & 2/8 cells & Universal claim not supported \\
    H4 & Adaptive risk dominates fixed controls & 3/12 comparisons & Universal dominance not supported \\
    \bottomrule
\end{tabular}
\end{center}
\end{samepage}
